\documentclass[12pt]{article}


\begin{document}


\title{Recent results on the Bergman metric of K\"ahler manifolds}
\author{Bo-Yong Chen}
\date{}
\maketitle


1. Let $M$ be a complex $n-$dimensional manifold. Let ${\cal H}$
be the Hilbert space of holomorphic $n-$forms on $M$ such that
$|\int_Mf\wedge \bar f|<\infty $. Let $h_0,h_1,\cdots $ be a
complete orthonormal basis for ${\cal H}$. Then the $2n-$form defined on $%
M\times M$ given by $B_M=\sum_{j=0}^\infty h_j\wedge \bar h_j$ is
called the \emph{Bergman kernel form} of $M$. Let
$(z_1,\cdots,z_n)$ be a local coordinate system in $M$ and let
$B_M(z)=B_M^{*}(z)dz_1\wedge \cdots \wedge dz_n\wedge d\bar
z_1\wedge \cdots \wedge d\bar z_n$ where $B_M^{*}$ is a locally
defined function. Thus $\beta_M :=\partial \bar \partial \log
B_M^{*}$ is a well-defined
Hermitian form of bidegree (1,1), whenever $B_M^{*}$ is nonzero. We say that $%
M $ possesses a \emph{Bergman metric} iff $\beta_M $ is everywhere
positive definite. The most important fact is that $\beta_M$ is an
invariant metric.

\bigskip

\emph{Example.} Any complete K\"ahler manifold with a bounded
continuous strictly psh function carries a Bergman metric.

\bigskip

2. \textbf{Bergman completeness and hyperconvexity}

\bigskip

\textbf{Definition.} A complex manifold $M$ is called hyperconvex
if there is a smooth strictly psh function $\rho$ on $M$ such that
$\{x\in M:\rho(x)<c\}\subset\subset M$ for all $c<0$.

\bigskip

\textbf{Theorem} (Blocki-Pflug '98, Herbort '99). \emph{Any
bounded hyperconvex domain in ${\bf C}^n$ is complete with respect
to the Bergman metric}.

\bigskip

Our first result is to generalize the above theorem as follows.

\bigskip

\textbf{Theorem 1.} \emph{Any hyperconvex manifold possess a
complete Bergman metric}.

\bigskip

Let us see some examples of hyperconvex manifolds:

\bigskip

(1) (Ohsawa-Sibony 98') pseudoconvex domains with $C^2$ boundary
in ${\bf P}^n$;

(2) (Greene-Wu 79') any complete K\"ahler manifold $M$ with a pole
$o$ such that its sectional curvature ${\rm K}$ is non-positive
and in addition satisfies
$$
{\rm K}\le -\frac{1+\epsilon}{r^2\log r}
$$
for some constant $\epsilon>0$ outside a compact subset of $M$,
where $r$ denotes the distance function based at $o$;

(3) (V$\hat{a}$j$\hat{a}$itu '96) holomorphic fiber bundles whose
base and fiber are hyperconvex.

\bigskip

3. \textbf{Bergman completeness and holomorphic convexity}.

\bigskip

A well-known theorem of Bremermann ('55) states that any bounded
Bergman complete domain in ${\bf C}^n$ is holomorphically convex.
This fact leads S. Kobayashi ('62) to ask the following

\bigskip

\textbf{Question}: \emph{Is a complex manifold holomorphically
convex if it possess a complete Bergman metric}?

\bigskip

Our second result is to give a counterexample. Namely, we show

\bigskip

\textbf{Theorem 2.} \emph{There is a Bergman complete manifold
which is not holomorphically convex}.

\bigskip

4. \textbf{Bergman metric on pseudoconvex domains in ${\bf P}^n$}

\bigskip

The following localization principle for the Bergman kernel and
metric was known for a long time.

\bigskip

\textbf{Theorem} (Ohsawa '81, Diederich-Fornaess-Herbort '84).
\emph{Let $\Omega\subset\subset {\bf C}^n$ be pseudoconvex. Then
for any two neighborhoods $V\subset\subset U$ of a boundary point
$p$ there is a constant $C>0$ such that }
\begin{eqnarray*}
B_\Omega(x) & \ge & C \cdot B_{\Omega\cap U
}(x)\\
\beta_\Omega(x;X) & \ge & C\cdot \beta_{\Omega\cap U }(x;X)
\end{eqnarray*}
\emph{for all $x\in \Omega\cap V$ and $X\in T^{1,0}_x(\Omega)$}.

\bigskip

It is natural to ask on which pseudoconvex domains in ${\bf P}^n$
the localization principle holds (Notice that the complement of a
hypersurface in ${\bf P}^n$ is pseudoconvex but the Bergman kernel
form vanishes). In this direction, Diederich-Ohsawa ('98) proved
the localization principle for another Bergman kernel (metric)
which is induced by $L^2$ holomorphic functions with respect to
the Fubini-Study metric on any pseudoconvex domain in ${\bf P}^n$
such that the interior of its complement is not empty. Such a
Bergman metric is \emph{not} an invariant metric.

\bigskip

\textbf{Theorem 3.} \emph{Let $\Omega$ be a pseudoconvex domain
with $C^2$ boundary in ${\bf P}^n$. Then there is a constant
$\gamma>0$ such that for any two neighborhoods $V\subset\subset U$
of a boundary point $p$, one has }
\begin{eqnarray*}
|B_\Omega(x)|_{FS} & \ge & C\cdot\frac{|B_{\Omega\cap U
}(x)|_{FS}}{|\log\delta_\Omega (x)|^{\gamma}}\\
\beta_\Omega(x;X) & \ge & C\cdot\frac{\beta_{\Omega\cap U
}(x;X)}{|\log\delta_\Omega (x)|^{\gamma}}
\end{eqnarray*}
\emph{for all $x\in \Omega\cap V$ and $X\in T^{1,0}_x(\Omega)$.
Here $\delta_\Omega$ $($resp. $|\cdot|_{FS})$ denotes the boundary
distance $($resp. length$)$ w.r.t. the Fubini-Study metric
$ds_{FS}^2$}.

\bigskip

Combining Theorem 3 with the Ohsawa-Takegoshi extension theorem,
we obtain

\bigskip

\textbf{Corollary.}
$$
|B_\Omega(x)|_{FS}  \ge  C\cdot
\delta_\Omega(x)^{-2}|\log\delta_\Omega|^{-\gamma}.
$$

\bigskip

We mention that Theorem 3 can not be generalized to an arbitrary
fat pseudoconvex domain in ${\bf P}^n$ as the following example
shows.

\bigskip

\emph{Example.} We consider the following domain
$$
\Omega=\{[z_0,z_1,z_2]\in {\bf P}^2:|z_1|<|z_0|\}
$$
where $(z_0,z_1,z_2)$ denote homogeneous coordinates in ${\bf
P}^2$. Note that $\Omega$ is biholomorphic to the product of the
unit disc and the complex plane via
$$
\zeta_1=z_1/z_0,\ \zeta_2=z_2/z_0.
$$
Hence $\Omega$ is a fat pseudoconvex domain in ${\bf P}^n$ but the
Bergman metric is degenerate.

\bigskip

However, the estimate of the Bergman metric in Theorem 3 does not
imply the completeness. On the other hand, Diederich-Ohsawa ('95)
proved that the Bergman distance for a bounded $C^2$ pseudoconvex
domain in ${\bf C}^n$ has a lower bound of a constant multiple
$\log|\log\delta_\Omega|$. It was improved by Blocki (02') to
$|\log \delta_\Omega|/\log|\log \delta_\Omega|$.

\bigskip

\textbf{Theorem 4.} \emph{Let $\Omega\subset {\bf P}^n$ be a
pseudoconvex domain with $C^2$ boundary. Then the Bergman distance
satisfies the following estimate:}
\[
{\rm dist}_\Omega (x_0,x)\ge \frac{C|\log\delta_\Omega(x)|}{\log
|\log \delta_\Omega (x)|}.
\]

\bigskip

5. \textbf{Sketch of proofs of theorems}

\bigskip

S. Kobayashi ('59, '62) proved that a complex manifold is Bergman
complete if the following condition is verified:

\bigskip

{\bf KC}. {\em Let $M$ be a complex manifold which
possesses a Bergman metric. Assume that there exists a dense subspace ${\cal %
H}^{\prime }$ of ${\cal H}$ such that for any $f\in {\cal
H}^{\prime }$ and for any sequence of points $\{y_k\}_{k=1}^\infty
$ of $M$ which has no adherent point in $M$, there is a
subsequence $\{y_{k_j}\}_{j=1}^\infty $ such that}
\[
\ \lim_{j\rightarrow \infty }\frac{f(y_{k_j})\wedge \bar f(y_{k_j})}{%
B_M(y_{k_j})}=0.
\]


\bigskip

It should be pointed out that W. Zwonek ('02) gave an example
which is Bergman complete but does not satisfy KC.

\bigskip

Let $g_M(\cdot,y)$ be the \emph{pluricomplex Green function} on
$M$, i.e.,
\[ g_M(x,y)=\sup \{u(x)\}
\]
where the superum is taken over all negative functions $u\in
PSH(M)$ satisfying the property that the function $u-\log |z|$ is
bounded from above in a deleted
neighborhood of $y$ for some holomorphic local coordinates $z$ centered at $%
y $, that is, $z(y)=0$. It satisfies the following property:

(1) $F:M\rightarrow M'$ biholomorphic$\Longrightarrow
g_M(x,y)=g_M(F(x),F(y))$;

(2) $M$ Stein, $\Omega\subset\subset M$
hyperconvex$\Longrightarrow (dd^c g_\Omega(\cdot,y))^n=(2\pi)^n
\delta_y$.

\bigskip

\emph{Proof of Theorem 1}.

i) The following Chern-Levine-Nirenberg type inequality due to
Blocki will play an important role:

\bigskip

\textbf{Proposition 5}. \emph{ Suppose $M$ is a complex manifold
and $\Omega\subset\subset M$ has $C^\infty$ boundary. Assume that
$u,v\in C^\infty(\overline{\Omega})$ are non-positive psh
functions such that $u|_{\partial \Omega}=0$. Then }
$$
\int_\Omega |u|^n (dd^c v)^n\le n! \|v\|_\infty^{n-1} \int_\Omega
|v|(dd^c u)^n.
$$

\bigskip

ii) By the above proposition, one has
$$
\int_M |g_M(\cdot,y)|^n (dd^c\rho)^n\le n! (2\pi)^n |\rho(y)|.
$$

iii) Let $f\in ${\em ${\cal H}$ }and {\em $\{y_k\}_{k=1}^\infty $
}be a sequence of points{\em \ }which has no adherent point in
$M$. By using the Donnelly-Fefferman $L^2$ estimate in such a way
as in T. Ohsawa's works, one can show
$$
\frac{f(y_{k})\wedge \bar f(y_{k})}{B_M(y_{k})}\leq C \left|
\int_{\{g_M(\cdot,y_k)<-1\}}f\wedge
 \overline{f}\right|.
$$
To verify KC, it suffices to prove that the latter tends to zero
as $k\rightarrow \infty$.

iv) For any $\epsilon>0$, there exists $M_\epsilon\subset\subset M
$ so that
$$
\left|\int_{M\backslash M_\epsilon}f\wedge
\bar{f}\right|<\epsilon.
$$
Notice that
\begin{eqnarray*}
\int_{M_\epsilon\cap \{g_M(\cdot,y_k)<-1\}} (dd^c\rho)^n &\le &
\int_M
|g_M(\cdot,y_k)|^n (dd^c\rho)^n\\
&\le & n! (2\pi)^n |\rho(y_k)|.
\end{eqnarray*}
This shows that one can choose a $k_\epsilon$ such that for all
$k>k_\epsilon$,
$$
\left|\int_{M_\epsilon\cap \{g_M(\cdot,y_k)<-1\}} f\wedge \bar{f}
\right|\le \sup_{M_\epsilon}\left|\frac{f\wedge \bar{f}}{(dd^c
\rho )^n}\right|\cdot \int_{M_\epsilon\cap \{g_M(\cdot,y_k)<-1\}}
(dd^c\rho)^n<\epsilon.
$$

\bigskip

\emph{Construction of the counterexample}

\bigskip

Let  $(L,h)$ be a holomorphic Hermitian line bundle over a
K\"ahler manifold $M$. Consider the Grauert tube:
$$
T(L,h)=\{w\in L:|w|_h<1\}.
$$

Let $S$ be a holomorphic $(n+1)-$form on $T(L,h)$. Locally,
$S=f(z,w)dz_1\wedge\cdots\wedge dz_n\wedge dw$ with
$f(z,w)=\sum_{j=0}^\infty f_j(z)w^j$. For each $j$, $f_j$ can be
viewed as a local representation of some holomorphic section $s_j$
of the adjoint line bundle $K_M+(j+1)L^\ast$. Notice that the
intersection of $T(L,h)$ with a fiber is the disk
$|w|^2<e^{-\varphi}$ and
$$
\frac{\sqrt{-1}}{2(j+1)\pi}\int_{|w|^2<e^{-\varphi}}|w^j|^2
dw\wedge d\bar{w}=e^{-(j+1)\varphi}
$$
if we write $h=e^{\varphi}$ locally. Therefore, if $S$ is $L^2$ on
$T(L,h)$, then $s_j$ belongs to
\begin{eqnarray*}
&& H^2(M,K_M+(j+1)L^\ast)\\
& = & \left\{s\in H^0(M,K_M+(j+1)L^\ast):\int_M |s|^2_{\omega^\ast
h^{-j-1}}dV_\omega<+\infty\right\}
\end{eqnarray*}
and
\begin{equation}
\left|\int_{T(L,h)}S\wedge \bar{S}\right|=\sum_{j=0}^\infty
2(j+1)\pi\|s_j\|^2_{(j+1)}.
\end{equation}
Here $\omega^\ast$ is the Hermitian metric on $K_M$ induced by
$\omega$ and $\|\cdot\|_{(j+1)}$ denotes the $L^2$ norm of
$H^2(M,K_M+(j+1)L^\ast)$. Such an observation goes back to Grauert
('62). Hence

\bigskip

\textbf{Proposition 6.} \emph{Suppose}

(i) \emph{$K_M+L^\ast$ is very ample in $H^2(M,K_M+L^\ast)$};

(ii) \emph{$K_M+2L^\ast$ is base point free in
$H^2(M,K_M+2L^\ast)$.}

\emph{Then $T(L,h)$ carries a Bergman metric.}

\bigskip


(1) reduces the verification of KC to certain adjoint line bundle
convexity of the base manifold. For example, the following holds:

\bigskip

\textbf{Proposition 7.} \emph{Let $(M,\omega)$ be a complete
K\"ahler manifold with injectivity radius $\tau>0$ and sectional
curvature $\le b^2$. Suppose there is a negative line bundle
$\pi:L\rightarrow M$ such that the curvature $ i\Theta(h)\le -
c\,\omega$ where $c$ is a positive constant. Then there is a
positive integer $m_0$ such that $T(mL,h^m)$ carries a complete
Bergman metric for all $m\ge m_0$.}

\bigskip

\emph{Proof.} Let $p_k=(z_k,w_k)\in T(mL,h^m)$. For each $L^2$
holomorphic $(n+1)$ form $S$ on $T(mL,h^m)$, one can associate it
with a sequence of sections $s_j\in H^2(M,K_M+(j+1)L^\ast)$. By
Hessian comparison theorem and $L^2$ estimate, one can construct
another section $\tilde{s}_j$ for each $j$ such that
$\tilde{s}_j(z_k)=s_j(z_k)$ and
$$
\|\tilde{s}_j\|_{(j+1)}\le C \|s_j\|_{B(z_k,\tau)}
$$
provided $m$ sufficiently large. Here $C$ is independent of $j,k$.
The local representations of $\{\tilde{s}_j\}$ give a holomorphic
$(n+1)$ form $\tilde{S}$ on $T(mL,h^m)$ with
$\tilde{S}(p_k)=S(p_k)$ and
$$
\|\tilde{S}\|\le C'|\int_{\pi^{-1}(B(z_k,\tau))} S\wedge \bar{S}|,
$$
thus verifies KC since $B(p_k)\ge \tilde{S}\wedge
\bar{\tilde{S}}(p_k)/\|\tilde{S}\|^2$.

\bigskip

Now let $p:M\rightarrow X$ be a holomorphic covering over a
projective manifold $X$. It is easy to verify that
$M,p^\ast(\omega),p^\ast(L)$ satisfy the condition of Proposition
2 , where $\omega$ is a K\"ahler metric of $X$ and $L$ is a
negative line bundle over $X$.

The counterexample can be constructed by the following classical
example: consider the Abelian variety ${\bf C}^2/\Lambda$ where
$\Lambda$ is the lattice generated by the elements
$(1,0),\,(0,1),\,(a,0)$ and $(0,b)$ over ${\bf Z}$ for complex
numbers $a,\,b$. Let $\tilde{\Lambda}$ be the subgroup of
$\Lambda$ generated by $(1,0),\,(0,1),\,(a,b)$ over ${\bf Z}$.
Note that ${\bf C}^2/\tilde{\Lambda}$ is a covering of ${\bf
C}^2/\Lambda$ while the latter is biholomorphic to the product of
one dimensional complex tori. It is easy to show that every
holomorphic function on ${\bf C}^2/\tilde{\Lambda}$ must be
constant provided $1,\,a,\,b$ are linearly independent over ${\bf
Z}$. One only needs to take $M={\bf C}^2/\tilde{\Lambda}$.

\bigskip

A well-known theorem of Siu states that very Stein submanifold of
a complex manifold has a Stein neighborhood. In fact, one can say
more

\bigskip

\textbf{Theorem 8}. \emph{Every Stein submanifold of a complex
manifold has a Stein neighborhood which carries a complete Bergman
metric.}

\bigskip

\emph{Proof of Theorem 3}

\bigskip

The difficulty lies in that it is not known whether there is a
bounded strictly psh function in a neighborhood of $\partial
\Omega$. However, Ohsawa-Sibony ('98) showed that there exists a
smooth and strictly psh function $\rho:\Omega\rightarrow (-1,0)$
such that $|\rho|\approx \delta^\tau_\Omega$ for some $0<\tau< 1$
and
$$
\partial\bar{\partial}\rho\ge C |\rho| ds^2_{FS}.
$$
This function plays an essential role in the construction of
certain weight function.

\bigskip

i) Given $y\in \Omega$, set
\begin{eqnarray*}
\phi_y & = & \kappa(d_{FS}(x,y))\log d_{FS}(x,y)-1\\
 \varphi _y & =
& \chi \left( \log (-\log (-\rho ))-\log (-\log (-\rho
(y)))\right)\phi_y.
\end{eqnarray*}
Then $\varphi _y-\log (-\phi_y)-b\log (-\rho )$ is psh for
suitable large constant $b$.

\bigskip

ii) By using patching $\rho$ to the above psh function, one can
construct a candidate function for the extreme problem of
$g_\Omega$ so that
$$
g_\Omega(x,y)>-\frac{c\rho(x)}{\rho(y)}\cdot|\log(-\rho(y))|
$$
for any $x,y\in \Omega$ with $|\rho(x)|\le |\rho(y)|/2$.

\bigskip

iii) By a technique coming from the pluripotential theory due to
Blocki and Herbort, one has
\begin{eqnarray*}
&& \{g_\Omega(x,y)<-1\}\\
&\subset & \{C^{-1}|\rho(y)|\cdot|\log (-\rho(y))|^{-1}\le
|\rho(x)|\le C |\rho(y)|\cdot|\log(-\rho(y))|^n\}
\end{eqnarray*}
It should be mentioned that while using this technique, a
difficulty arise: unlike the case ${\bf C}^n$, there does not
necessary exist parallel transformations on manifolds. That is the
reason why the theorem can not be generalized to complex manifolds
with positive holomorphic bisectional curvature.

\bigskip

iv) Set
$\tilde{\epsilon}=C^{-1}|\rho(y)|\cdot|\log(-\rho(y))|^{-1}$ and
$\Omega_{\tilde{\epsilon}}=\{x\in
\Omega:|\rho(x)|>\tilde{\epsilon}\}$. Set
\begin{eqnarray*}
\tilde{\lambda}_y &= & \varphi_y-\log(-\phi_y)-b\log(-\rho)+b\log
\tilde{\epsilon}\\
\psi_y &= & \max\{\tilde{\lambda}_y,g_\Omega(\cdot,y)\}.
\end{eqnarray*}
$\psi_y$ a negative psh function $\psi_y$ with a log pole at $y$
such that
$$
\psi_y(x) \ge   -c-b\log|\log(-\rho(y))|, \ \forall\, x\in
\{\psi_y<-1\}\backslash B(y,r)
$$
where $r,b,c$ are independent of $y$ and $B(y,r)$ denotes the
geodesic ball.

\bigskip

vi) Use the following $L^2$ estimate

\bigskip

{\bf Theorem} (Diederich-Ohsawa '95, Berndtsson-Charpentier '00).
{\em Let $M$ be a Stein manifold. Let $\varphi ,\psi $ be psh
functions such that $r\partial \bar
\partial \psi \ge
\partial \psi \bar \partial \psi $ holds in the sense of distribution for some $0<r<1$. Then for any
$\bar \partial-$closed $(n,1)$ form $v$ with $\int_M|v|_{\partial
\bar
\partial (\varphi +\psi )}^2e^{\psi -\varphi}<+\infty$, there exists an $(n,0)$ form $u$ on
$M$ so that $\bar
\partial u=v$ and}
\[
\left| \int_Mu\wedge \bar ue^{\psi -\varphi }\right| \le
C_r\int_M|v|_{\partial \bar \partial (\varphi +\psi )}^2e^{\psi
-\varphi }.
\]

\bigskip

We apply this theorem with
$\psi=-\frac12\log(-\rho),\,\varphi=2n\psi_y-\log(-\psi_y+1)+\psi$
to solve
$$
\bar{\partial}u=\bar{\partial}(\lambda \left( -\log
(-\psi_y+1)+\log 2\right) \kappa (d_{FS}(\cdot,y)/2r)f),
$$
and obtain an $F\in {\cal H}(\Omega_{\tilde{\epsilon}})$ with
$F(y)=f(y)$ and
$$
\|F\|\le C(\log\delta_\Omega)^\gamma \|f\|.
$$

\bigskip


vii) Since $\{g_\Omega(\cdot,y)<-1\}\subset
\Omega_{\tilde{\epsilon}}$, by using the above $L^2$ estimate
again, we extend $F$ to $\Omega$.

\bigskip

6. \textbf{An open problem}

\bigskip

A domain $\Omega$ in a complex manifold is called locally
hyperconvex if for each $p\in \Omega$ there is a neighborhood $U$
of $p$ so that $\Omega\cap U$ is hyperconvex.

Diederich and Ohsawa ('98) constructed a locally hyperconvex
domain in ${\bf P}^n$ which is not hyperconvex.

\bigskip

\textbf{Question.} \emph{Does a locally hyperconvex domain in
${\bf P}^n$ carry a complete Bergman metric}?


\begin{thebibliography}{99}
\bibitem{2}  B. Berndtsson and Ph. Charpentier, {\em A Sobolev mapping
property of the Bergman kernel}, Math. Z. {\bf 235} (2000), 1--10.

\bibitem{}  Z. Blocki, {\em Estimates for the complex Monge-Ampere operator}%
, Bull. Pol. Acad. Sci. {\bf 41} (1993), 151--157.

\bibitem{}  ----------, \emph{The Bergman metric and the pluricomplex Green
function}, MPI (Leipzig) preprint no: 85 (2002).

\bibitem{1}  Z. Blocki and P. Pflug, {\em Hyperconvexity and Bergman
completeness}, Nagoya Math. J. {\bf 151} (1998), 221--225.

\bibitem{} B. Y. Chen, \emph{Bergman completeness of hyperconvex
manifolds}, preprint.

\bibitem{} ----------, \emph{On certain K\"ahler manifolds which possess a Bergman
metric}, preprint.

\bibitem{} ----------, \emph{On Bergman metric of pseudoconvex domains in a complex projective space}, preprint.

\bibitem{}  K. Diederich and T. Ohsawa, {\em An estimate for the Bergman
distance on pseudoconvex domains}, Ann. of Math. {\bf 141} (1995),
181--190.

\bibitem{} ----------, \emph{On pseudoconvex domains in ${\bf
P}^n$}, Tokyo J. Math. {\bf 21} (1998), 353--358.

\bibitem{10}  G. Herbort, {\em The Bergman metric on hyperconvex domains},
Math. Z. {\bf 232} (1999), 183--196.


\bibitem{}  S. Kobayashi, {\em On complete Bergman metrics}, Proc. Amer. Math.
Soc. {\bf 13} (1962), 511--513.

\bibitem{} T. Ohsawa and N. Sibony, \emph{Bounded p.s.h. functions and
pseudoconvexity in K\"ahler manifolds}, Nagoya Math. J. {\bf 149}
(1998), 1--8.
\end{thebibliography}
\end{document}